\section{The conditional landing chain}
\label{sec:landing_chain}

This section proves the conditional theorems and the concrete
calibration.  The reader may keep the bare shell in mind throughout: for
it \(c(\rho)=\Real\rho\), the represented zeros are all nontrivial
zeros, and \Cref{thm:rh:conditional,thm:rh:full-classical} say the same
thing.  The duality-confinement theorem is used in the abstract
form of the companion paper \citep[Remark~7.2]{TsiokosSDTC2026}: if
\(a\) is a positive element of an ordered set, \(\tau\) is a
nonnegative monotone real trace with \(\tau(a)=0\) only if \(a=0\), and
\(a\preceq b_n\) with \(\tau(b_n)\to0\), then \(a=0\), because
\(0\leq\tau(a)\leq\tau(b_n)\to0\).

\begin{theorem}[Duality-confinement theorem applied to \(\Sel\)]
\label{thm:rh:dc-master-applied}
Let \(\Sel\) be a shell for which \(\GamSDTCSelberg\) holds
(\Cref{obl:rh:gamma-sdtc-selberg}).  Then the
anti-invariant ledger of the represented zeros vanishes: \(\AZ=0\).
\end{theorem}

\begin{proof}
By items (ii) and (iii) and the abstract squeeze, \(\AX\) is the zero
element; by item (ii\('\)), \(\psim(x)=0\) at every visible object \(x\).
For a represented zero \(\rho\), the object \(x_\rho\) is visible, so
by item (iv) \(e(c(\rho)-\frac12)=\psim(x_\rho)=0\), and
hence \(c(\rho)=\frac12\).  (Separation, item (i), gives the equivalent
statement that \(x_\rho\) is fixed by \(J\); it is not needed for the
conclusion.)  Every summand of \(\AZ\) therefore vanishes, and
\(\AZ=0\) by \Cref{thm:rh:translation-T-reverse} applied to the represented zeros.
\end{proof}

The companion theorem enters this proof only through the abstract
squeeze; the link from \(\AX\) to the readout is item (ii\('\)) and the
link to the zeros is item (iv), both of which are assumed.

\begin{theorem}[Conditional landing chain]
\label{thm:rh:conditional}
Let \(\Sel\) be a shell for which \(\GamSDTCSelberg\) holds.  Then every represented zero \(\rho\) of \(\Sel\)
satisfies \(c(\rho)=\frac12\).  On the bare shell this says
\(\Real\rho=\frac12\) for every nontrivial zero \(\rho\).
\end{theorem}

\begin{proof}
\Cref{thm:rh:dc-master-applied} gives \(\AZ=0\), and the forward
direction of Theorem~T (\Cref{thm:rh:translation-T-forward}), applied
to the represented zeros, gives the conclusion.
\end{proof}

\begin{theorem}[Full classical conditional Riemann hypothesis]
\label{thm:rh:full-classical}
Let \(\Sel\) be a shell with coordinate system \(R\), and suppose that
\(\GamSDTCSelberg\) holds for \(\Sel\) and that there is a
\emph{classical-zero identification}: a map \(t:R\to\mathbb R\) with \(t(\frac12)=\frac12\)
such that for every nontrivial zero \(s\) of \(\zeta\) there is a
represented zero \(\rho\) with \(t(c(\rho))=\Real s\).
Then the Riemann hypothesis holds.  For the bare shell, \(t\) is the
identity.
\end{theorem}

\begin{proof}
By \Cref{thm:rh:conditional}, \(c(\rho)=\frac12\) for every
represented zero, so \(\Real s=t(\frac12)=\frac12\) for every nontrivial
zero \(s\).  The Riemann hypothesis follows, because every zero of \(\zeta\) other
than the trivial zeros lies in the strip \(0<\Real s<1\) (proof of
\Cref{thm:rh:translation-T}).
\end{proof}

Taken together with \Cref{prop:rh:gamma-calibration}, these theorems
say that the Riemann hypothesis is equivalent to the statement that \(\GamSDTCSelberg\) holds for the bare shell.  The recognition source
used there is degenerate: a single object carries the readout.  The
next theorem applies the operator form of the duality-confinement
theorem to the actual zeros, with the zeros themselves as objects.
  Let
\((W_n)_{n\geq1}\) be finite sets of nontrivial zeros, each closed under
\(\JL\), such that every nontrivial zero lies in \(W_n\) for all large
\(n\); such a sequence exists because the nontrivial zeros are countable and
each orbit of \(\JL\) has one or two elements (take \(W_n\) to be the
union of the orbits of the first \(n\) zeros in an enumeration).  On
\(W_n\) take the finite weighted ledger of the companion
paper with unit weights, response space \(\mathbb R\), response
involution \(v\mapsto-v\), involution \(\JL\) and readout
\(\psimRH(\rho)=\Real\rho-\frac12\).  Its operator is multiplication by the unit-weight window sum
\(A_{W_n}=\sum_{\rho\in W_n}(\Real\rho-\frac12)^2\), and its readout is
separating: it vanishes exactly at the fixed points of \(\JL\).  The
response space is one-dimensional, so the operators are scalars and the
theorem is the simplest case of the duality-confinement theorem; what
it adds to \Cref{prop:rh:gamma-calibration} is that the objects are the
actual zeros and the bridge is an identity rather than an
assumption.

\begin{theorem}[Concrete window form]
\label{thm:rh:concrete-window-dc}
With the windows \(W_n\) above, the following are equivalent:
\begin{enumerate}[label=\textup{(\alph*)}]
  \item for every \(n\) there are real numbers \(b_{n,k}\geq A_{W_n}\)
    with \(b_{n,k}\to0\) as \(k\to\infty\);
  \item there are real numbers \(b_n\geq A_{W_n}\) with \(b_n\to0\);
  \item the Riemann hypothesis.
\end{enumerate}
In (a) and (b) the numbers \(b\) are the traces of positive operators
on \(\mathbb R\) dominating the window operators, so (a) is the
hypothesis of the duality-confinement theorem for each window, and the
implication (a)\(\Rightarrow\)(c) is proved by that theorem.
\end{theorem}

\begin{proof}
(a)\(\Rightarrow\)(c): for each \(n\), \(0\leq A_{W_n}\leq b_{n,k}\to0\), so
\(A_{W_n}=0\); this is the duality-confinement theorem for this finite
weighted ledger \citep[Theorem~7.1]{TsiokosSDTC2026}, and since the readout
is separating every zero in \(W_n\) is fixed by \(\JL\), that is, lies on
the critical line; every nontrivial zero lies in some
window, and \Cref{thm:rh:translation-T} gives the Riemann hypothesis.
(b)\(\Rightarrow\)(c): if some nontrivial zero \(\rho_0\) were off the
line, it would lie in \(W_n\) for all large \(n\), so
\(b_n\geq A_{W_n}\geq(\Real\rho_0-\frac12)^2>0\) for all large \(n\),
contradicting \(b_n\to0\).  (c)\(\Rightarrow\)(a) and
(c)\(\Rightarrow\)(b): under the Riemann hypothesis every window sum is
zero, and the zero numbers dominate it.
\end{proof}

Unwound, condition (a) says that every window sum is zero and
condition (b) that the window sums tend to zero; the theorem makes
explicit that, for the actual zeros, domination records of vanishing
trace are nothing more than such upper bounds, and that they exist
exactly when the window sums vanish.  A
proof of the Riemann hypothesis along this route therefore needs an
independent source of such bounds; \Cref{sec:finite_windows} examines
some candidate sources.
